\section*{Capítulo \ref{ch:b2:unicellular-diversity} --- Diversidad de los organismos unicelulares}

\begin{solution}{exo:b2:unicellular-diversity:1}
\omterm{def:b2:unicellular-diversity:unicellular}{Unicelular}: una célula realiza todas las funciones y vive sola;
\omterm{def:b2:unicellular-diversity:unicellular}{colonial}: células idénticas permanecen unidas, pero cada una podría vivir
sola; \omterm{def:b2:unicellular-diversity:unicellular}{pluricelular}: varios tipos celulares, interdependientes, y solo la
línea germinal se reproduce. \emph{E.~coli} y la levadura son
\omterm{def:b2:unicellular-diversity:unicellular}{unicelulares} (la yema se desprende y vive sola); \emph{Gonium} es
\omterm{def:b2:unicellular-diversity:unicellular}{colonial} (16 células idénticas, cada una capaz de fundar una colonia);
\emph{Volvox} es \omterm{def:b2:unicellular-diversity:unicellular}{pluricelular} (células somáticas mortales, unas pocas
células germinales).
\end{solution}

\begin{solution}{exo:b2:unicellular-diversity:2}
$S/V = 3/r$: \qty{6}{\micro m^{-1}} para el coco,
\qty{0.06}{\micro m^{-1}} para el \omterm{def:b2:unicellular-diversity:protist}{protista}, una diferencia de cien veces.
El coco respira más deprisa por unidad de volumen: cualquier punto de su
citoplasma está a una fracción de micrómetro de la superficie por la que
entran el oxígeno y el alimento, mientras que el interior del \omterm{def:b2:unicellular-diversity:protist}{protista} se
abastece a través de cien veces menos superficie por unidad de volumen.
\end{solution}

\begin{solution}{exo:b2:unicellular-diversity:3}
Sin núcleo (un \omterm{def:b2:unicellular-diversity:prokaryote}{nucleoide}) frente a una envoltura nuclear; sin
mitocondrias ni endomembranas, con la cadena respiratoria en la membrana
plasmática, frente a orgánulos; un flagelo que es una hélice rígida de
flagelina girada por protones, frente a un cilio $9+2$ que se curva
gracias a la dineína y al ATP; además, el tamaño (\qty{1}{\micro m}
frente a decenas) y una pared de peptidoglucano. Una bacteria no puede
fagocitar: su pared le impide englobar una partícula, algo que una ameba
hace con sus pseudópodos.
\end{solution}

\begin{solution}{exo:b2:unicellular-diversity:4}
Cilios: locomoción y barrido del alimento hacia la boca; surco oral:
canaliza las partículas hacia la citofaringe; vacuola digestiva: digiere
las bacterias englobadas con enzimas lisosómicas; \omterm{prop:b2:unicellular-diversity:osmoregulation}{vacuola contráctil}:
expulsa el agua que entra por ósmosis; macronúcleo: el núcleo funcional,
con cientos de copias de cada gen transcritas para la vida diaria;
micronúcleo: el núcleo diploide germinal, que interviene en la meiosis y
la conjugación.
\end{solution}

\begin{solution}{exo:b2:unicellular-diversity:5}
$t = x^2/2D$: $(10^{-6})^2/10^{-9} = \qty{1}{ms}$ para la bacteria;
$(5\times 10^{-4})^2/10^{-9} = \qty{250}{s}$, unos cuatro minutos, para
la ameba. La ameba no puede confiar en la difusión para repartir los
metabolitos: necesita corrientes citoplasmáticas y un transporte
impulsado por motores a lo largo del citoesqueleto, de los que la
bacteria prescinde.
\end{solution}

\begin{solution}{exo:b2:unicellular-diversity:6}
$Re = \rho vL/\eta$: \emph{Paramecium} $10^3\times 10^{-3}\times
2\times 10^{-4}/10^{-3} = 0.2$; bacteria $10^3\times 2\times
10^{-5}\times 2\times 10^{-6}/10^{-3} = 4\times 10^{-5}$; nadador
$10^3\times 1\times 2/10^{-3} = 2\times 10^6$. A $Re \gg 1$ la inercia
lleva al nadador más allá tras la brazada; a $Re \ll 1$ el arrastre
viscoso detiene a la bacteria en menos de un nanómetro, de modo que solo
avanza mientras bate.
\end{solution}

\begin{solution}{exo:b2:unicellular-diversity:7}
$v_s = 2r^2\Delta\rho g/9\eta = 2\times 10^{-10}\times 100\times
9.81/(9\times 10^{-3}) = \qty{2.2e-5}{m/s}$; descender \qty{50}{m} lleva
$50/2.2\times 10^{-5} = \qty{2.3e6}{s}$, 27 días. Reducir el radio a la
mitad divide $v_s$ por 4: 108 días. Reducir $\Delta\rho$ a
\qty{20}{kg/m^3} la divide por 5: 135 días.
\end{solution}

\begin{solution}{exo:b2:unicellular-diversity:8}
$V = \tfrac{4}{3}\pi\times 100\times 25\times 25 =
\qty{2.6e5}{\micro m^3}$. En \qty{1800}{s}: $2.6\times 10^5/1800 =
\qty{145}{\micro m^3/s}$; como $\qty{1}{\micro m^3} =
\qty{1e-15}{L}$, esto equivale a \qty{1.5e-13}{L/s}.
\end{solution}

\begin{solution}{exo:b2:unicellular-diversity:9}
Una duplicación a lo largo de \qty{1}{mm} es un aumento de alrededor del
\qty{0.07}{\%} por micrómetro (puesto que $2^{1/1000} - 1 = 0.0007$), es
decir, un \qty{0.07}{\%} de un extremo a otro de la célula. Una carrera de
\qty{1}{s} recorre \qty{20}{\micro m}: \qty{1.4}{\%}. Para resolver una
diferencia del \qty{0.07}{\%} la célula tendría que contar unas
$(1/0.0007)^2 = 2\times 10^6$ moléculas en cada extremo, muchas más de
las que se unen en un segundo a sus pocos miles de receptores; el
\qty{1.4}{\%} requiere unas \num{5000}, lo que es factible. La
comparación temporal convierte una medida espacial imposible en una
posible, dejando que la carrera haga la integración.
\end{solution}

\begin{solution}{exo:b2:unicellular-diversity:10}
Superficie $4\pi r^2$ con $r = \qty{250}{\micro m}$: \qty{7.9e5}{\micro
m^2}; volumen citoplasmático $\approx$ superficie $\times$ \qty{2}{\micro m}
$= \qty{1.6e6}{\micro m^3}$; relación \qty{0.5}{\micro m^{-1}}. Para
\emph{E.~coli}, un cilindro: $S/V = 2\pi r L/\pi r^2 L = 2/r =
\qty{4}{\micro m^{-1}}$. La bacteria gigante solo está ocho veces peor
que \emph{E.~coli}, no 500 veces, porque todo punto de su citoplasma
está a menos de \qty{2}{\micro m} de la superficie. La vacuola almacena
nitrato, su aceptor de electrones, para poder respirar sulfuro en el
sedimento anóxico entre las escasas agitaciones que traen nitrato
(\cref{ch:b2:microbial-metabolism}).
\end{solution}

\begin{solution}{exo:b2:unicellular-diversity:11}
Una célula no puede nadar y dividirse a la vez. Una colonia de 16
células que deja de nadar durante las horas de la división se hunde,
según la \omterm{thm:b2:unicellular-diversity:stokes}{ley de Stokes}, una distancia despreciable (su radio es de unas
decenas de micrómetros y su velocidad de hundimiento, de micrómetros por
segundo); toda la colonia puede dividirse sin perder nada. Una esfera de
\num{10000} células y \qty{500}{\micro m} de radio se hundiría cien
veces más deprisa --- milímetros por segundo, metros en una hora ---
fuera de la luz: compensa mantener nadando a la mayoría de las células
mientras unas pocas se dividen. Las células germinales son grandes
porque deben llevar los recursos para construir, mediante divisiones
sucesivas y sin alimentarse, una esfera hija entera de miles de células;
por eso las abastece el soma.
\end{solution}

\begin{solution}{exo:b2:unicellular-diversity:12}
Un clado es un antepasado y todos sus descendientes; un grado es un
nivel de organización alcanzado por varios linajes. Los \omterm{def:b2:unicellular-diversity:protist}{protistas} ---
algas verdes (junto a las plantas terrestres), diatomeas (junto a las
algas pardas), ciliados (junto a los dinoflagelados y \emph{Plasmodium}),
amebas (cerca de los animales y los hongos) --- solo comparten la célula
eucariota que vive sola, y su último antepasado común es el de todos los
eucariotas, nosotros incluidos: un grado. ``\omterm{def:b2:unicellular-diversity:prokaryote}{Procariota}'' reúne bacterias
y arqueas por lo que les falta; las arqueas están más cerca de los
eucariotas que de las bacterias, así que el término vuelve a designar un
grado. Las cianobacterias descienden de un único antepasado que inventó
la fotosíntesis oxigénica e incluyen a todos sus descendientes (dejando
aparte el linaje de los cloroplastos): un clado.
\end{solution}

\begin{solution}{pb:b2:unicellular-diversity:1}
\textbf{1.} $1\times 1\times 0.1 = \qty{0.1}{mm^3} = \qty{0.1}{\micro L}$.
\textbf{2.} $24/\qty{0.1}{\micro L} = 240$ por \unit{\micro L} $=
\qty{2.4e5}{}$ células por mililitro.
\textbf{3.} $V = \tfrac{4}{3}\pi\times 125 = \qty{524}{\micro m^3}$;
$S/V = 3/5 = \qty{0.6}{\micro m^{-1}}$.
\textbf{4.} $2.4\times 10^5\times 524 = \qty{1.26e8}{\micro m^3}$ por
mililitro, y $\qty{1}{mL} = \qty{1e12}{\micro m^3}$: una fracción de
$1.3\times 10^{-4}$, el \qty{0.013}{\%}.
\textbf{5.} $\qty{1000}{m^3} = \qty{1e9}{mL}$: $2.4\times 10^{14}$
células.
\textbf{6.} La cianobacteria es un filamento de células mucho más
pequeñas (unos pocos micrómetros), verdeazuladas, sin cloroplasto ni
núcleo diferenciados, y no nada con flagelos; el alga es una sola célula
ovoide con dos flagelos, un estigma y un único cloroplasto en forma de
copa.
\textbf{7.} $384/24 = 16 = 2^4$: cuatro duplicaciones en \qty{48}{h},
$T = \qty{12}{h}$.
\textbf{8.} $k = \ln 2/T = 0.693/12 = \qty{0.058}{h^{-1}}$.
\textbf{9.} $t = \ln(10^8/2.4\times 10^5)/k = \ln(417)/0.058 =
6.03/0.058 = \qty{104}{h}$, unos 4,3 días.
\textbf{10.} Se agotan los nutrientes (fosfato, nitrato); las células se
sombrean unas a otras, de modo que la luz por célula disminuye; los
consumidores (ciliados, rotíferos, \emph{Daphnia}) se multiplican a su
costa; a \qty{1e8}{} por mililitro las células llenarían alrededor del
cinco por ciento del volumen y agotarían el $\mathrm{CO_2}$
disuelto.
\textbf{11.} Crecimiento solo la mitad del tiempo: la tasa media se
reduce a la mitad, \qty{208}{h}, 8,7 días.
\textbf{12.} El recuento mide la población, no células individuales: si
una célula libera 8 hijas cada \qty{36}{h}, la población se multiplica
por 8 $= 2^3$ en tres \omterm{def:b2:unicellular-diversity:division}{tiempos de generación} de \qty{12}{h} cada uno. El
\omterm{def:b2:unicellular-diversity:division}{tiempo de generación} es el tiempo medio de duplicación del número de
células, cualquiera que sea la manera en que se agrupen las divisiones.
\textbf{13.} $Re = 10^3\times 10^{-4}\times 10^{-5}/10^{-3} = 10^{-3}$.
\textbf{14.} $m = 1050\times 5.24\times 10^{-16} = \qty{5.5e-13}{kg}$;
$6\pi\eta r = 6\pi\times 10^{-3}\times 5\times 10^{-6} =
\qty{9.4e-8}{kg/s}$; $d = 5.5\times 10^{-13}\times 10^{-4}/9.4\times
10^{-8} = \qty{5.8e-10}{m}$: medio nanómetro.
\textbf{15.} $v_s = 2\times 25\times 10^{-12}\times 50\times
9.81/(9\times 10^{-3}) = \qty{2.7e-6}{m/s}$, \qty{2.7}{\micro m/s}.
\textbf{16.} Hundirse \qty{1}{m}: $1/2.7\times 10^{-6} =
\qty{3.7e5}{s}$, 4,3 días; subir nadando a \qty{100}{\micro m/s}:
\qty{1e4}{s}, 2,8 horas. Nadar gana por un factor de 37.
\textbf{17.} $c_\infty = \qty{1e-2}{mol/m^3}$: $J = 4\pi\times
1.9\times 10^{-9}\times 5\times 10^{-6}\times 10^{-2} =
\qty{1.2e-15}{mol/s} = 7.2\times 10^8$ moléculas por segundo.
\textbf{18.} \qty{100}{pg} $= \qty{1e-10}{g} = \qty{8.3e-12}{mol}$ de
carbono, duplicado en \qty{12}{h} $= \qty{43200}{s}$:
\qty{1.9e-16}{mol/s} $= 1.2\times 10^8$ átomos por segundo. Aporte
entre demanda: $7.2/1.2 = 6$.
\textbf{19.} El aporte se multiplica por diez ($\propto r$) y la demanda
por mil ($\propto r^3$): la razón pasa a ser $6/100 = 0.06$; una célula
así solo podría crecer al seis por ciento del ritmo, de modo que un alga
de ese tamaño alimentada por difusión es imposible sin agitación,
transporte interno o un metabolismo mucho más lento.
\textbf{20.} $\Delta\Pi = RT\Delta c = 8.314\times 293\times 99 =
\qty{2.4e5}{Pa}$, \qty{2.4}{bar}.
\textbf{21.} Superficie $4\pi(5\times 10^{-6})^2 = \qty{3.14e-10}{m^2}$;
entrada $L_p A\Delta\Pi = 10^{-14}\times 3.14\times 10^{-10}\times
2.4\times 10^5 = \qty{7.6e-19}{m^3/s} = \qty{0.76}{\micro m^3/s}$.
\textbf{22.} $524/0.76 = \qty{690}{s}$, unos once minutos.
\textbf{23.} $P = \Delta\Pi\times$ caudal $= 2.4\times 10^5\times
7.6\times 10^{-19} = \qty{1.8e-13}{W}$.
\textbf{24.} Un ATP: $5\times 10^4/6.02\times 10^{23} =
\qty{8.3e-20}{J}$: $1.8\times 10^{-13}/8.3\times 10^{-20} = 2.2\times
10^6$ ATP por segundo. Fijación del carbono: $3\times 1.2\times 10^8 =
3.6\times 10^8$ ATP por segundo. La vacuola cuesta el \qty{0.6}{\%} de
esa cifra --- un seguro barato contra el estallido.
\textbf{25.} Entrada de agua de \qty{0.76}{\micro m^3/s} (el volumen de
la célula en once minutos); sin \omterm{prop:b2:unicellular-diversity:osmoregulation}{vacuola contráctil} la célula duplicaría
su volumen en unos \qty{700}{s}; mantenerse a salvo del agua cuesta al
menos $2\times 10^6$ ATP por segundo, menos del uno por ciento de lo que
su fotosíntesis gasta en fijar carbono.
\end{solution}
